% \iffalse meta-comment
% !TEX program = pdfLaTeX
%<*internal>
\iffalse
%</internal>
%<*readme>
================================================================

Working with the command \fbox or \fcolorbox, one has to
handle page breaks by hand. The present package defines the
environment mdframed which automatically deals with page breaks.

Author's name: Marco Daniel and Elke Schubert (!new)
License type: lppl

================================================================
%</readme>
%<*internal>
\fi
\def\nameofplainTeX{plain}
\ifx\fmtname\nameofplainTeX\else
  \expandafter\begingroup
\fi
%</internal>
%<*install>
\input docstrip.tex
\keepsilent
\askforoverwritefalse
\preamble
----------------------------------------------------------------
Working with the command fbox or fcolorbox, one has to
handle page breaks by hand. The present package defines the
environment mdframed which automatically deals with page breaks.

Author's name: Marco Daniel and Elke Schubert (!new)
License type: lppl
\endpreamble
\postamble

================================================================
Copyright (C) 2011 by Marco Daniel

This work may be distributed and/or modified under the
conditions of the LaTeX Project Public License (LPPL), either
version 1.3c of this license or (at your option) any later
version.  The latest version of this license is in the file:

http://www.latex-project.org/lppl.txt

This work is "maintained" (as per LPPL maintenance status) by
Marco Daniel.

Have fun!

================================================================
\endpostamble
\usedir{doc/latex/mdframed}
\generate{\file{mdframed-example-default.tex}{\from{mdframed-examples.dtx}{mdframed-example-default}}}
\generate{\file{mdframed-example-tikz.tex}{\from{mdframed-examples.dtx}{mdframed-example-tikz}}}
\generate{\file{mdframed-example-pstricks.tex}{\from{mdframed-examples.dtx}{mdframed-example-pstricks}}}
\generate{\file{mdframed-example-texsx.tex}{\from{mdframed-examples.dtx}{mdframed-example-texsx}}}
\Msg{*********************************************************}
\Msg{*}
\Msg{* To finish the installation you have to move the}
\Msg{* following file into a directory searched by TeX:}
\Msg{*}
\Msg{* \space\space documentation.sty}
\Msg{*}
\Msg{* To produce the documentation run the file documentation.dtx}
\Msg{* once through LaTeX. Then, run}
\Msg{*}
\Msg{* \space\space makeindex -s gglo.ist -o mdframed.gls mdframed.glo}
\Msg{* \space\space makeindex -s gind.ist mdframed.idx}
\Msg{*}
\Msg{* through makeIndex to produce the glossary. Finally, run LaTeX once again.}
\Msg{* That's all!}
\Msg{*}
\Msg{* Happy TeXing!}
\Msg{*********************************************************}
%</install>
%<install>\endbatchfile
%<*internal>
\ifx\fmtname\nameofplainTeX
  \expandafter\endbatchfile
\else
  \expandafter\endgroup
\fi
%</internal>
%<*driver>
%%$Id: mdframed-examples.dtx 256 2011-12-04 18:32:06Z marco $
\documentclass[parskip=false,11pt,]{ltxmdf}
\usepackage{mdframed}
\ltxmdfsetifoot $Id: mdframed-examples.dtx 256 2011-12-04 18:32:06Z marco $
\CodelineIndex
\RecordChanges
\begin{document}
  \DocInput{\jobname.dtx}
\end{document}
%</driver>
% \fi
%
% \CheckSum{725}
%
% \CharacterTable
%  {Upper-case    \A\B\C\D\E\F\G\H\I\J\K\L\M\N\O\P\Q\R\S\T\U\V\W\X\Y\Z
%   Lower-case    \a\b\c\d\e\f\g\h\i\j\k\l\m\n\o\p\q\r\s\t\u\v\w\x\y\z
%   Digits        \0\1\2\3\4\5\6\7\8\9
%   Exclamation   \!     Double quote  \"     Hash (number) \#
%   Dollar        \$     Percent       \%     Ampersand     \&
%   Acute accent  \'     Left paren    \(     Right paren   \)
%   Asterisk      \*     Plus          \+     Comma         \,
%   Minus         \-     Point         \.     Solidus       \/
%   Colon         \:     Semicolon     \;     Less than     \<
%   Equals        \=     Greater than  \>     Question mark \?
%   Commercial at \@     Left bracket  \[     Backslash     \\
%   Right bracket \]     Circumflex    \^     Underscore    \_
%   Grave accent  \`     Left brace    \{     Vertical bar  \|
%   Right brace   \}     Tilde         \~}
%
% \GetFileInfo{mdframed-examples}
%
%
%
% \title{Examples of the \Pack{mdframed} package}
% \subtitle{auto-split frame environment}
% \author{\href{mailto:marco.daniel@mada-nada.de}{Marco Daniel}%
%        \footnote{With thanks to Heiko Oberdiek, Rolf Niepraschk,
%                  Martin Scharrer and Herbert Voss.\newline Sorry for bad English.}
%           \href{mailto:marco.daniel@mada-nada.de}{Elke Schubert}}
% \version{\mdversion}
% \date{\today}
% \introduction{The standard methods for framing text (\Cmd{fbox} or \Cmd{fcolorbox}) require you to handle page breaks by hand, meaning that you have to split the \Cmd{fbox} into two. The present package defines the environment \Pack{mdframed} which automatically deals with pagebreaks in framed text.\par
% By defining new environments the user may choose between several individual designs.%
% \par\kern.5\baselineskip\noindent
% The dtx-file creates only the files \begin{tabular}[t]{ll}
%             \texttt{mdframed-example-default.tex} &
%             \texttt{mdframed-example-tikz.tex} \\
%             \texttt{mdframed-example-pstricks.tex}  &
%             \texttt{mdframed-example-texsx.tex}
%             \end{tabular}
% \par\kern.5\baselineskip\noindent
% FYI:\quad\parbox[t]{.8\linewidth}{I create a repository for \Pack{mdframed} on \href{https://github.com/marcodaniel/mdframed}{github} where you can \href{https://github.com/marcodaniel/mdframed}{download} the current development status.}
% }
% 
% 
% 
% 
% \maketitle
% 
% 
% \section{Motivation}
% To complie the files use |Makefile| with |make EXAMPLES|.
%\StopEventually{^^A
%  \PrintChanges^^A
%}
%\iffalse
%    \begin{macrocode}
%<*mdframed-example-default>
%    \end{macrocode}
%\fi
%
% \subsection{The file mdframed-example-default}

%    \begin{macrocode}
%Documenation of the package mdframed
%%$Id: mdframed-examples.dtx 256 2011-12-04 18:32:06Z marco $
\setcounter{errorcontextlines}{999}
\documentclass[parskip=false,english,11pt]{ltxmdf}
\ltxmdfsetifoot  $Id: mdframed-examples.dtx 256 2011-12-04 18:32:06Z marco $

\usepackage{showexpl}
\lstset{style=lstltxmdf,explpreset={pos=b,rframe={}},}

\newcommand\Loadedframemethod{default}
\usepackage[framemethod=\Loadedframemethod]{mdframed}

\title{The \Pack{mdframed} package}
\subtitle{Examples for \Opt{framemethod=\Loadedframemethod}}
\author{\href{mailto:marco.daniel@mada-nada.de}{Marco Daniel}}
\version{\mdversion}
\date{\today}
\introduction{In this document I collect various examples for \Opt{framemethod=\Loadedframemethod}.
Some presented examples are more or less exorbitant.}

\mdfsetup{skipabove=\topskip,skipbelow=\topskip}
\newrobustcmd\ExampleText{%
        An \textit{inhomogeneous linear} differential equation has the form
         \begin{align}
            L[v ] = f,
         \end{align}
        where $L$ is a linear differential operator, $v$ is
        the dependent variable, and $f$ is a given non-zero
        function of the independent variables alone.
}

\newcounter{examplecount}
\setcounter{examplecount}{0}
\renewcommand\thesubsection{}
\newcommand\Examplesec[1]{%
\stepcounter{examplecount}%
\subsection{Example~\arabic{examplecount}~--~#1\relax}%
}

\begin{document}
\maketitle
\section{Loading}
In the preamble only the package \Pack{mdframed} width the option \Opt{framemethod=\Loadedframemethod} is loaded. All other modifications will be done by \Cmd{mdfdefinestyle} or \Cmd{mdfsetup}.

{\large\color{red!50!black}
\NOTE Every \Cmd{global} inside the examples is necessary to work with the package \Pack{showexpl}.}

\section{Examples}
All examples have the following settings:

\begin{tltxmdfexample}
\mdfsetup{skipabove=\topskip,skipbelow=\topskip}
\newrobustcmd\ExampleText{%
An \textit{inhomogeneous linear} differential equation
has the form
\begin{align}
L[v ] = f,
\end{align}
where $L$ is a linear differential operator, $v$ is
the dependent variable, and $f$ is a given non-zero
function of the independent variables alone.
}
\end{tltxmdfexample}
\clearpage
\Examplesec{very simple}
\begin{LTXexample}
\global\mdfdefinestyle{exampledefault}{%
     linecolor=red,linewidth=3pt,%
     leftmargin=1cm,rightmargin=1cm
}
\begin{mdframed}[style=exampledefault]
\ExampleText
\end{mdframed}
\end{LTXexample}

\Examplesec{hidden line + frame title}
\begin{LTXexample}
\global\mdfapptodefinestyle{exampledefault}{%
 topline=false,rightline=true,bottomline=false}
\begin{mdframed}[style=exampledefault,frametitle={Inhomogeneous linear}]
\ExampleText
\end{mdframed}
\end{LTXexample}
\clearpage

\Examplesec{colored frame title}
\begin{LTXexample}
\renewcommand\mdframedtitleenv[1]{%
                 \colorbox{green}{%
                 \parbox{\dimexpr\linewidth-6pt\relax}%6pt=linewidth
                 {\centering\bfseries #1}}%
                 \par\kern.5\baselineskip\noindent%
        }
\global\mdfapptodefinestyle{exampledefault}{%
   rightline=true,innerleftmargin=0,innerrightmargin=0}
\begin{mdframed}[style=exampledefault,frametitle={Inhomogeneous linear}]
\ExampleText
\end{mdframed}
\end{LTXexample}

\Examplesec{framed picture which is centered}
\begin{LTXexample}
\begin{mdframed}[userdefinedwidth=6cm,align=center,
                 linecolor=blue,linewidth=4pt]
\includegraphics[width=\linewidth]{donald-duck}
\end{mdframed}
\end{LTXexample}

\Examplesec{theorem with separate header and the help of TikZ (complex)}
\begin{tltxmdfexample}
\makeatletter
\newcounter{theo}[section]
\newcommand*\newmdframedtitleenv[1]{%
 \@afterindentfalse
  {\parindent \z@
    \setlength{\parfillskip}{\z@ plus 1fil}%
    \mdraggedtitle\nobreak%
    \makebox[\linewidth][l]{%
     \hspace*{-1\mdf@innerleftmargin@length}%
     \rlap{\color{white}%
       \hspace*{-1\mdf@middlelinewidth@length}%
       \rule[\mdf@middlelinewidth@length]%
            {\dimexpr\linewidth+1\mdf@innerleftmargin@length%
             +\mdf@innerrightmargin@length
             +2\mdf@middlelinewidth@length\relax}%
            {\dimexpr\ht\strutbox+.3333em\relax}%
     }%
     \rlap{\color{blue!20}%
       \rule{\dimexpr\linewidth+\mdf@innerleftmargin@length%
             +\mdf@innerrightmargin@length\relax}%
            {\mdf@middlelinewidth@length}}%
     \hspace*{-1\mdf@middlelinewidth@length}%
     \tikz[remember picture,baseline]%
       \node[,draw = none, text = black,fill = blue!20,]%
       {\mdf@frametitlefont\strut Theorem~\thetheo#1};\relax%
   }%
  \par\kern.5\baselineskip}%
 \@afterheading}
\newenvironment{theo}[1][]{%
  \let\mdframedtitleenv\newmdframedtitleenv%
   \stepcounter{theo}%
   \ifstrempty{#1}%
     {\mdfsetup{frametitle={\strut}}}%
     {\mdfsetup{frametitle={:~#1}}}%
   \begin{mdframed}[innertopmargin=0pt,linecolor=blue!20,%
                    linewidth=2pt,topline=false,]%
   }{\end{mdframed}}
\begin{theo}[Inhomogeneous Linear]
\ExampleText
\end{theo}

\begin{theo}
\ExampleText
\end{theo}
\end{tltxmdfexample}
\makeatletter
\newcounter{theo}[section]
\newcommand*\newmdframedtitleenv[1]{%
 \@afterindentfalse
  {\parindent \z@
    \setlength{\parfillskip}{\z@ plus 1fil}%
    \mdraggedtitle\nobreak%
    \makebox[\linewidth][l]{%
     \hspace*{-1\mdf@innerleftmargin@length}%
     \rlap{\color{white}%
       \hspace*{-1\mdf@middlelinewidth@length}%
       \rule[\mdf@middlelinewidth@length]%
            {\dimexpr\linewidth+1\mdf@innerleftmargin@length%
             +\mdf@innerrightmargin@length
             +2\mdf@middlelinewidth@length\relax}%
            {\dimexpr\ht\strutbox+.3333em\relax}%
     }%
     \rlap{\color{blue!20}%
       \rule{\dimexpr\linewidth+\mdf@innerleftmargin@length%
             +\mdf@innerrightmargin@length\relax}%
            {\mdf@middlelinewidth@length}}%
     \hspace*{-1\mdf@middlelinewidth@length}%
     \tikz[remember picture,baseline]%
       \node[,draw = none, text = black,fill = blue!20,]%
       {\mdf@frametitlefont\strut Theorem~\thetheo#1};\relax%
   }%
  \par\kern.5\baselineskip}%
 \@afterheading}
\newenvironment{theo}[1][]{%
  \let\mdframedtitleenv\newmdframedtitleenv%
   \stepcounter{theo}%
   \ifstrempty{#1}%
     {\mdfsetup{frametitle={\strut}}}%
     {\mdfsetup{frametitle={:~#1}}}%
   \begin{mdframed}[innertopmargin=0pt,linecolor=blue!20,%
                    linewidth=2pt,topline=false,]%
   }{\end{mdframed}}

\begin{theo}[Inhomogeneous Linear]
\ExampleText
\end{theo}

\begin{theo}
\ExampleText
\end{theo}

\clearpage
\Examplesec{hide only a part of a line}
The example below is inspired by the following post on StackExchange \href{http://tex.stackexchange.com/questions/24101/theorem-decorations-that-stay-with-theorem-environment}{Theorem decorations that stay with theorem environment}
\begin{LTXexample}
\makeatletter
\newlength{\interruptlength}
\setlength{\interruptlength}{2.5ex}
\newrobustcmd\overlaplines{%
 \appto\md@frame@leftline@single{%
  \llap{\color{white}%
   \rule[\dimexpr-\mdfboundingboxdepth%
         \ifbool{mdf@bottomline}{-\mdf@middlelinewidth@length}{}%
         +\interruptlength\relax]%
       {\mdf@middlelinewidth@length}%
       {\dimexpr\mdfboundingboxtotalheight%
        +\ifbool{mdf@bottomline}{\mdf@middlelinewidth@length}{0pt}
        +\ifbool{mdf@topline}{\mdf@middlelinewidth@length}{0pt}%
       -2\interruptlength\relax}%
  }%
 }%
 \appto\md@frame@rightline@single{%
  \rlap{\color{white}%
    \hspace*{\mdfboundingboxwidth}%
    \hspace*{\mdf@innerrightmargin@length}%
    \rule[\dimexpr-\mdfboundingboxdepth%
          \ifbool{mdf@bottomline}{-\mdf@middlelinewidth@length}{}
          +\interruptlength\relax]%
         {\mdf@middlelinewidth@length}%
         {\dimexpr\mdfboundingboxtotalheight%
          +\ifbool{mdf@bottomline}{\mdf@middlelinewidth@length}{0pt}%
          +\ifbool{mdf@topline}{\mdf@middlelinewidth@length}{0pt}
          -2\interruptlength\relax}%
  }%
 }
}
\makeatother
\overlaplines

\begin{mdframed}[linecolor=blue,linewidth=8pt]
\ExampleText
\end{mdframed}
\end{LTXexample}
\end{document}
 \endinput
%    \end{macrocode}
%
%\iffalse
%    \begin{macrocode}
%</mdframed-example-default>
%    \end{macrocode}
%\fi
%
%\iffalse
%    \begin{macrocode}
%<*mdframed-example-tikz>
%    \end{macrocode}
%\fi
%
% \subsection{The file mdframed-example-tikz}

%    \begin{macrocode}
%Documenation of the package mdframed
%%$Id: mdframed-examples.dtx 256 2011-12-04 18:32:06Z marco $
\setcounter{errorcontextlines}{999}
\documentclass[parskip=false,english,11pt]{ltxmdf}
\ltxmdfsetifoot  $Id: mdframed-examples.dtx 256 2011-12-04 18:32:06Z marco $

\usepackage{showexpl}
\lstset{style=lstltxmdf,explpreset={pos=b,rframe={}},}

\newcommand\Loadedframemethod{TikZ}
\usepackage[framemethod=\Loadedframemethod]{mdframed}

\title{The \Pack{mdframed} package}
\subtitle{Examples for \Opt{framemethod=\Loadedframemethod}}
\author{\href{mailto:marco.daniel@mada-nada.de}{Marco Daniel}}
\version{\mdversion}
\date{\today}
\introduction{In this document I collect various examples for \Opt{framemethod=\Loadedframemethod}.
Some presented examples are more or less exorbitant.}

\mdfsetup{skipabove=\topskip,skipbelow=\topskip}
\newrobustcmd\ExampleText{%
        An \textit{inhomogeneous linear} differential equation has the form
         \begin{align}
            L[v ] = f,
         \end{align}
        where $L$ is a linear differential operator, $v$ is
        the dependent variable, and $f$ is a given non-zero
        function of the independent variables alone.
}

\newcounter{examplecount}
\setcounter{examplecount}{0}
\renewcommand\thesubsection{}
\newcommand\Examplesec[1]{%
\stepcounter{examplecount}%
\subsection{Example~\arabic{examplecount}~--~#1\relax}%
}

\begin{document}
\maketitle
\section{Loading}
In the preamble only the package \Pack{mdframed} width the option \Opt{framemethod=\Loadedframemethod} is loaded. All other modifications will be done by \Cmd{mdfdefinestyle} or \Cmd{mdfsetup}.

{\large\color{red!50!black}
\NOTE Every \Cmd{global} inside the examples is necessary to work with the package \Pack{showexpl}.}

\section{Examples}
All examples have the following settings:

\begin{tltxmdfexample}
\mdfsetup{skipabove=\topskip,skipbelow=\topskip}
\newrobustcmd\ExampleText{%
An \textit{inhomogeneous linear} differential equation
has the form
\begin{align}
L[v ] = f,
\end{align}
where $L$ is a linear differential operator, $v$ is
the dependent variable, and $f$ is a given non-zero
function of the independent variables alone.
}
\end{tltxmdfexample}
\clearpage
\ExampleText{round corner}
\begin{LTXexample}
\global\mdfdefinestyle{exampledefault}{%
     outerlinewidth=5pt,innerlinewidth=0pt,
     outerlinecolor=red,roundcorner=5pt
}
\begin{mdframed}[style=exampledefault]
\ExampleText
\end{mdframed}
\end{LTXexample}

\Examplesec{hidden line + frame title}
\begin{LTXexample}
\global\mdfapptodefinestyle{exampledefault}{%
 topline=false,leftline=false,}
\begin{mdframed}[style=exampledefault,frametitle={Inhomogeneous linear}]
\ExampleText
\end{mdframed}
\end{LTXexample}
\clearpage
\Examplesec{framed picture which is centered}
\begin{LTXexample}
\begin{mdframed}[userdefinedwidth=6cm,align=center,
                 linecolor=blue,middlelinewidth=4pt,roundcorner=5pt]
\includegraphics[width=\linewidth]{donald-duck}
\end{mdframed}
\end{LTXexample}

\Examplesec{Gimmick}
\begin{LTXexample}
\mdfsetup{splitbottomskip=0.8cm,splittopskip=0cm,
          innerrightmargin=2cm,innertopmargin=1cm,%
          innerlinewidth=2pt,outerlinewidth=2pt,
          middlelinewidth=10pt,backgroundcolor=red,
          linecolor=blue,middlelinecolor=gray,
          tikzsetting={draw=yellow,line width=3pt,%
                    dashed,%
                    dash pattern= on 10pt off 3pt},
          rightline=false,bottomline=false}
\begin{mdframed}
\ExampleText
\end{mdframed}
\end{LTXexample}

\Examplesec{complex example with TikZ}

\begin{tltxmdfexample}
\tikzstyle{titregris} =
          [draw=gray, thick, fill=white, shading = exersicetitle, %
           text=gray, rectangle, rounded corners,
           right,minimum height=.7cm]

\pgfdeclarehorizontalshading{exersicebackground}{100bp}
{color(0bp)=(green!40);
color(100bp)=(black!5)}

\pgfdeclarehorizontalshading{exersicetitle}{100bp}
{color(0bp)=(red!40);
color(100bp)=(black!5)}

\newcounter{exercise}
\renewcommand\theexercise{Exercise~n\arabic{exercise}}
\makeatletter
\def\mdf@@exercisepoints{}
\define@key{mdf}{exercisepoints}{%
    \def\mdf@@exercisepoints{#1}
}
\renewrobustcmd\mdfcreateextratikz{%
      \node[titregris,xshift=1cm] at (P-|O) %
           {~\mdf@frametitlefont{\theexercise}~};
      \ifdefempty{\mdf@@exercisepoints}%
      {}%
      {\node[titregris,left,xshift=-1cm] at (P)%
        {~\mdf@frametitlefont{\mdf@@exercisepoints points}~};}%
}
\makeatother

\mdfdefinestyle{exercisestyle}{%
  outerlinewidth=1pt,
  innerlinewidth=0pt,
  roundcorner=2pt,
  linecolor=gray,
  tikzsetting={shading = exersicebackground},
  innertopmargin=1.2\baselineskip,
  skipabove={\dimexpr0.5\baselineskip+\topskip\relax},
  needspace=3\baselineskip,
  frametitlefont=\sffamily\bfseries,
  settings={\global\stepcounter{exercise}},
  }

\begin{mdframed}[style=exercisestyle,]
\ExampleText
\end{mdframed}

\begin{mdframed}[style=exercisestyle,exercisepoints=10]
\ExampleText
\end{mdframed}
\end{tltxmdfexample}

\tikzstyle{titregris} =
          [draw=gray, thick, fill=white, shading = exersicetitle, %
           text=gray, rectangle, rounded corners,
           right,minimum height=.7cm]

\pgfdeclarehorizontalshading{exersicebackground}{100bp}
{color(0bp)=(green!40);
color(100bp)=(black!5)}

\pgfdeclarehorizontalshading{exersicetitle}{100bp}
{color(0bp)=(red!40);
color(100bp)=(black!5)}

\newcounter{exercise}
\renewcommand\theexercise{Exercise~n\arabic{exercise}}
\makeatletter
\def\mdf@@exercisepoints{}
\define@key{mdf}{exercisepoints}{%
    \def\mdf@@exercisepoints{#1}
}
\renewrobustcmd\mdfcreateextratikz{%
      \node[titregris,xshift=1cm] at (P-|O) {~\textbf{\theexercise}~};
      \ifdefempty{\mdf@@exercisepoints}%
      {}%
      {\node[titregris,left,xshift=-1cm] at (P)%
        {~\mdf@frametitlefont{\mdf@@exercisepoints points}~};}%
}
\makeatother

\mdfdefinestyle{exercisestyle}{%
  outerlinewidth=1pt,
  innerlinewidth=0pt,
  roundcorner=2pt,
  linecolor=gray,
  tikzsetting={shading = exersicebackground},
  innertopmargin=1.2\baselineskip,
  skipabove={\dimexpr0.5\baselineskip+\topskip\relax},
  needspace=3\baselineskip,
  frametitlefont=\sffamily\bfseries,
  settings={\global\stepcounter{exercise}},
  }

\begin{mdframed}[style=exercisestyle,]
\ExampleText
\end{mdframed}

\begin{mdframed}[style=exercisestyle,exercisepoints=10]
\ExampleText
\end{mdframed}

\end{document}
 \endinput
%    \end{macrocode}
%
%\iffalse
%    \begin{macrocode}
%</mdframed-example-tikz>
%    \end{macrocode}
%\fi
%
%\iffalse
%    \begin{macrocode}
%<*mdframed-example-pstricks>
%    \end{macrocode}
%\fi
%
% \subsection{The file mdframed-example-pstricks}

%    \begin{macrocode}
%Documenation of the package mdframed
%%$Id: mdframed-examples.dtx 256 2011-12-04 18:32:06Z marco $
\setcounter{errorcontextlines}{999}
\documentclass[parskip=false,english,11pt]{ltxmdf}
\ltxmdfsetifoot$Id: mdframed-examples.dtx 256 2011-12-04 18:32:06Z marco $



\lstDeleteShortInline{|}
\newcommand\Loadedframemethod{PSTricks}
\usepackage[framemethod=\Loadedframemethod]{mdframed}

\usepackage{showexpl}
\lstset{style=lstltxmdf,explpreset={pos=b,rframe={}},}



\title{The \Pack{mdframed} package}
\subtitle{Examples for \Opt{framemethod=\Loadedframemethod}}
\author{\href{mailto:marco.daniel@mada-nada.de}{Marco Daniel}}
\version{\mdversion}
\date{\today}
\introduction{In this document I collect various examples for \Opt{framemethod=\Loadedframemethod}.
Some presented examples are more or less exorbitant.}

\mdfsetup{skipabove=\topskip,skipbelow=\topskip}
\newrobustcmd\ExampleText{%
        An \textit{inhomogeneous linear} differential equation has the form
         \begin{align}
            L[v ] = f,
         \end{align}
        where $L$ is a linear differential operator, $v$ is
        the dependent variable, and $f$ is a given non-zero
        function of the independent variables alone.
}

\newcounter{examplecount}
\setcounter{examplecount}{0}
\renewcommand\thesubsection{}
\newcommand\Examplesec[1]{%
\stepcounter{examplecount}%
\subsection{Example~\arabic{examplecount}~--~#1\relax}%
}

\begin{document}
\maketitle
\section{Loading}
In the preamble only the package \Pack{mdframed} width the option \Opt{framemethod=\Loadedframemethod} is loaded. All other modifications will be done by \Cmd{mdfdefinestyle} or \Cmd{mdfsetup}.

{\large\color{red!50!black}
\NOTE Every \Cmd{global} inside the examples is necessary to work with the package \Pack{showexpl}.}
X
\section{Examples}
All examples have the following settings:

\begin{tltxmdfexample}
\mdfsetup{skipabove=\topskip,skipbelow=\topskip}
\newrobustcmd\ExampleText{%
An \textit{inhomogeneous linear} differential equation
has the form
\begin{align}
L[v ] = f,
\end{align}
where $L$ is a linear differential operator, $v$ is
the dependent variable, and $f$ is a given non-zero
function of the independent variables alone.
}
\end{tltxmdfexample}
\clearpage

\Examplesec{very simple}
\begin{LTXexample}
\global\mdfdefinestyle{exampledefault}{%
     linecolor=red,linewidth=3pt,%
     leftmargin=1cm,rightmargin=1cm
}
\begin{mdframed}[style=exampledefault,roundcorner=5]
\ExampleText
\end{mdframed}
\end{LTXexample}

\Examplesec{hidden line + frame title}
\begin{LTXexample}
\global\mdfapptodefinestyle{exampledefault}{%
 topline=false,rightline=false,bottomline=false}
\begin{mdframed}[style=exampledefault,frametitle={Inhomogeneous linear}]
\ExampleText
\end{mdframed}
\end{LTXexample}

\clearpage

\Examplesec{Dash Lines}
\begin{LTXexample}
\global\mdfdefinestyle{exampledefault}{%
   pstrickssetting={linestyle=dashed,},linecolor=red,linewidth=5pt}
\begin{mdframed}[style=exampledefault,]
\ExampleText
\end{mdframed}
\end{LTXexample}

\Examplesec{Double Lines}
\begin{LTXexample}
\global\mdfdefinestyle{exampledefault}{%
   pstrickssetting={doubleline=true,doublesep=3pt},
   linecolor=red,linewidth=5pt}
\begin{mdframed}[style=exampledefault,]
\ExampleText
\end{mdframed}
\end{LTXexample}
\end{document}
 \endinput
%    \end{macrocode}
%
%\iffalse
%    \begin{macrocode}
%</mdframed-example-pstricks>
%    \end{macrocode}
%\fi
%
%\iffalse
%    \begin{macrocode}
%<*mdframed-example-texsx>
%    \end{macrocode}
%\fi
%
% \subsection{The file mdframed-example-texsx}

%    \begin{macrocode}
%Documenation of the package mdframed
%%$Id: mdframed-examples.dtx 256 2011-12-04 18:32:06Z marco $
\setcounter{errorcontextlines}{999}
\documentclass[parskip=false,english,11pt,ltxlipsum]{ltxmdf}
\ltxmdfsetifoot  $Id: mdframed-examples.dtx 256 2011-12-04 18:32:06Z marco $

\usepackage{showexpl}
\lstset{style=lstltxmdf,explpreset={pos=b,rframe={}},}

\newcommand\Loadedframemethod{default}
\usepackage[framemethod=\Loadedframemethod]{mdframed}

\title{The \Pack{mdframed} package}
\subtitle{Examples for \Opt{framemethod=\Loadedframemethod}}
\author{\href{mailto:marco.daniel@mada-nada.de}{Marco Daniel}}
\version{\mdversion}
\date{\today}
\introduction{In this document I collect various examples for \Opt{framemethod=\Loadedframemethod}.
Some presented examples are more or less exorbitant.}

\mdfsetup{skipabove=\topskip,skipbelow=\topskip}
\newrobustcmd\ExampleText{%
        An \textit{inhomogeneous linear} differential equation has the form
         \begin{align}
            L[v ] = f,
         \end{align}
        where $L$ is a linear differential operator, $v$ is
        the dependent variable, and $f$ is a given non-zero
        function of the independent variables alone.
}

\newcounter{examplecount}
\setcounter{examplecount}{0}
\renewcommand\thesubsection{}
\newcommand\Examplesec[1]{%
\stepcounter{examplecount}%
\subsection{Example~\arabic{examplecount}~--~#1\relax}%
}

\begin{document}
\maketitle
\section{Loading}
In the preamble only the package \Pack{mdframed} width the option \Opt{framemethod=\Loadedframemethod} is loaded. All other modifications will be done by \Cmd{mdfdefinestyle} or \Cmd{mdfsetup}.

{\large\color{red!50!black}
\NOTE Every \Cmd{global} inside the examples is necessary to work with the package \Pack{showexpl}.}

\section{Examples}
All examples have the following settings:

\begin{tltxmdfexample}
\mdfsetup{skipabove=\topskip,skipbelow=\topskip}
\newrobustcmd\ExampleText{%
An \textit{inhomogeneous linear} differential equation
has the form
\begin{align}
L[v ] = f,
\end{align}
where $L$ is a linear differential operator, $v$ is
the dependent variable, and $f$ is a given non-zero
function of the independent variables alone.
}
\end{tltxmdfexample}
\clearpage
\Examplesec{Package listings}
The example below is inspired by the following post on StackExchange \href{http://tex.stackexchange.com/questions/27673/background-overflows-when-using-rounded-corners-for-listings-package-listings}{Background overflows when using rounded corners for listings (package: `listings`)}

Here the solution which can be decorate as usual.

\begin{tltxmdfexample}[moretexcs={BeforeBeginEnvironment,AfterEndEnvironment},morekeywords={lstlisting}]
\BeforeBeginEnvironment{lstlisting}{%
    \begin{mdframed}[<modification>]%
    \vspace{-0.7em}}
\AfterEndEnvironment{lstlisting}{%
    \vspace{-0.5em}%
    \end{mdframed}}
\end{tltxmdfexample}

\Examplesec{Package multicol}
How I wrote in \enquote{Known Problems} you can't combine \Pack{multicol}  with \Pack{mdframed}. In a simple way without any breaks you can use:
\begin{LTXexample}
\begin{multicols}{2}
\lipsum[1]
\begin{mdframed}
\ExampleText
\end{mdframed}
\lipsum[2]
\end{multicols}
\end{LTXexample}
\clearpage
\twocolumn[\Examplesec{Working in twocolumn mode}]
\begin{tltxmdfexample}
\twocolumn[%
  \Examplesec{Working in
          twocolumn mode}]
\lipsum[2]
\begin{mdframed}[%
   leftmargin=10pt,%
   rightmargin=10pt,%
   linecolor=red,
   backgroundcolor=yellow]
\ExampleText
\end{mdframed}
\lipsum[2]
\end{tltxmdfexample}
\lipsum[2]\lipsum[2]
\begin{mdframed}[leftmargin=10pt,%
                 rightmargin=10pt,%
                 linecolor=red,
                 backgroundcolor=yellow]
\ExampleText
\end{mdframed}
\lipsum[2]
\onecolumn
\Examplesec{Working inside enumerate}
\begin{LTXexample}
Text Text Text Text Text Text Text Text
\begin{enumerate}
\item in the following \ldots
      \begin{mdframed}[linecolor=blue,linewidth=2]
         \ExampleText
      \end{mdframed}
\item \lipsum[2]
\end{enumerate}
Text Text Text Text Text Text
\end{LTXexample}
\end{document}
 \endinput
%    \end{macrocode}
%
%\iffalse
%    \begin{macrocode}
%</mdframed-example-texsx>
%    \end{macrocode}
%\fi
%
% \Finale
%
\endinput

